\documentclass[11pt]{article}
\usepackage[a4paper,margin=28mm]{geometry}
\usepackage{amsmath,amssymb,amsthm,booktabs}
\usepackage[hidelinks]{hyperref}
\usepackage{xurl,microtype}
\setlength{\emergencystretch}{3em}
\newcommand{\Q}{\mathbb Q}
\newtheorem{proposition}{Proposition}
\title{A reciprocal chart saving one coordinate\\in the shared-reference quartic construction}
\author{ZQ Diophantine Project}
\date{23 November 2026}
\begin{document}
\maketitle
\begin{abstract}
We remove one coordinate and one quadratic equation from each of the two
general shared-reference constructions. The resulting complete systems
have 2085 coordinates and 1775 quadratic rows, or 2277 coordinates and 1943
quadratic rows. The change is a rational bijection using existing nonzero
units, and all six affected equations are displayed and checked exactly.
The numerical counts and finite algebra are verified. Any undecidability
interpretation remains conditional on the explicitly unresolved source
universality hypothesis \(U_{\mathrm{em}}\) and on the retained arithmetic
interfaces; this note asserts no unconditional undecidability bound.
\end{abstract}

The starting systems, their simultaneous completeness proof, and their
coefficient tables are given in the companion
\nolinkurl{shared_quartic_checkpoint.tex}. The exact source obligation
\(U_{\mathrm{em}}\) and the discrepancy between the consulted published
displays are explained there and in
\nolinkurl{jones1980_quartic_transfer.tex}. The present change is independent
of that obligation: it proves equivalence of the displayed rational
systems. The 56-coordinate ordinary-\(abc\) route is not altered.

\section{A reciprocal input without an additional coordinate}
\label{sec:reciprocal-reference-rechart}

The preceding shared-reference compilation introduced an ordinary input
\(i\) and a row \(iv=1\), where \(v\) is the outer coordinate denoted
\texttt{i\_inverse}.  This permits all reference functions to be shared
in their actual physical scale.  For the general 163- and 179-coordinate
interfaces, one can instead store \(i\) in an existing witness coordinate.
The following improvement is a rational change of coordinates, so it adds
no arithmetic assumption.  Its use as an undecidability reduction still
requires the unresolved hypothesis \(U_{\mathrm{em}}\) and every retained
arithmetic interface.

\begin{proposition}\label{prop:reciprocal-reference-rechart}
The shared-reference systems with \(2086\) and \(2278\) coordinates have
rationally equivalent presentations with, respectively,
\[
\begin{array}{c|cc}
\text{integrality interface}&\text{coordinates}&\text{quadratic rows}\\\hline
163&2085&1775\\
179&2277&1943.
\end{array}
\]
The outer input projection and the ten common reference functions are
preserved.  The sum of the squares of the new rows has degree exactly four.
\end{proposition}

\begin{proof}
Work in the integer block that certifies \(i=1/v\).  Denote its old
coordinates by \(Z_j\), with \(Z_1=i\).  In the 163-coordinate block let
\[
       q=Z_{159},\qquad q^{-1}=Z_{160},\qquad q^{-2}=Z_{161};
\]
in the 179-coordinate block, after the already recorded three-register
rechart, the corresponding indices are \(162,163,164\).  The two literal
unit rows force \(q\ne0\) and give the indicated inverse values.  The
outer equation \(vE=c^2\), with \(c>0\), already forces \(v\ne0\).

Replace three witnesses, retaining their coordinate slots, by \(t,h,r\).
For the 163-coordinate interface define the old witnesses from the new by
\begin{equation}\label{eq:reciprocal-chart-163}
       Z_{37}=\frac{t}{vq},\qquad
       Z_{77}=\frac{h}{vq},\qquad
       Z_{78}=\frac{r}{v}.
\end{equation}
For the 179-coordinate interface instead use
\begin{equation}\label{eq:reciprocal-chart-179}
       Z_{37}=\frac{t}{vq},\qquad
       Z_{77}=\frac{h}{vq},\qquad
       Z_{78}=\frac{r}{v^2q^2}.
\end{equation}
All other coordinates are unchanged.  These transformations and their
inverses are rational because \(v,q\ne0\).  None of the ten reference
coordinates has index \(37,77,78\), so their common physical values are
unaffected.

Here is the complete transformed subsystem.  Put \(R=Z_{90}\), and let
\((j_0,j_1)=(117,118)\) for the 163-coordinate interface and
\((j_0,j_1)=(118,119)\) for the 179-coordinate interface.  Four rows are
common to both presentations:
\begin{align}
 1-vt&=0,\label{eq:reciprocal-input-unit}\\
 2601t^2+255102tv-4hv+6255001v^2&=0,\label{eq:reciprocal-pair}\\
 v(-Z_{17}+Z_{23}+2504Z_{j_0})+51tZ_{j_0}&=0,
       \label{eq:reciprocal-cleared-real}\\
 v(-Z_{18}+Z_{24}+2504Z_{j_1})+51tZ_{j_1}&=0.
       \label{eq:reciprocal-cleared-imag}
\end{align}
For the 163-coordinate interface the remaining two rows are
\begin{align}
 15531948th+291790992tv-512h^2+767410608hv
       +14315517724v^2+vr&=0,\label{eq:reciprocal-product-163}\\
 22201rR+vR-149C^2r&=0,\label{eq:reciprocal-forced-163}
\end{align}
where \(C\) is its coefficient parameter, not a quantified coordinate.
For the fixed 179-coordinate interface they are
\begin{align}
 15531948th+291790992tv-512h^2+767410608hv
       +14315517724v^2+rR&=0,\label{eq:reciprocal-product-179}\\
 25rR-A_*r+v^2&=0.\label{eq:reciprocal-forced-179}
\end{align}
Here \(A_*\) is exactly the fixed integer in that interface's original
row \nolinkurl{scalar.forced_prime.definition}; no coefficient is changed.
Its decimal expansion is stored, without abbreviation, in the generator
and in each complete coefficient table.

For clarity, every row multiplication can be checked directly.  First set
\(Z_1=1/v\) in the old equations.  The old input-square equation is
\(Z_1^2-qZ_{37}=0\); multiplication by \(v^2\) gives
\eqref{eq:reciprocal-input-unit}.  The old
\nolinkurl{scalar.rPair1.definition} and
\nolinkurl{scalar.four_contact_product.definition} rows, after the
appropriate substitution, become the displayed pair and product rows
when multiplied by \(v^2q^2\).  The forced-prime row is multiplied by
\(v\) in the 163-coordinate case and by \(v^2q^2\) in the
179-coordinate case.  Each old \texttt{input.cleared} row is multiplied
by \(vq\).  These multipliers are all nonzero.  An occurrence audit of
the complete source table shows that these six rows are the only rows
containing \(Z_1,Z_{37},Z_{77},Z_{78}\).  Thus no unlisted equation
requires alteration, and every new row has degree at most two.

On a solution of the old shared-reference system, define the new witnesses
by the inverse of \eqref{eq:reciprocal-chart-163} or
\eqref{eq:reciprocal-chart-179}.  Since \(iv=1\), the old input-square
row gives \(t=vqi^2/q=vi^2=i\).  Hence the new system is satisfied after
deleting the old external coordinate \(i\) and its row \(iv=1\).
Conversely, a solution of the new system satisfies \(vt=1\).  Reconstruct
\(i=t=1/v\), reconstruct the three old witnesses using the displayed
formulas, and divide the six new rows by their nonzero multipliers.  This
recovers every old equation, including \(iv=1\).  The two constructions
are inverse on rational solutions and fix all outer coordinates.

One coordinate and one row have therefore been removed in each case.
The new input-unit row has the nonzero quadratic term \(-vt\).
Over \(\mathbb Q\), the sum of the squares of the quadratic homogeneous
parts cannot vanish identically: evaluating at a real point with a
nonzero such part gives a positive value.  Consequently the sum of the
squares of all residuals has degree exactly four and vanishes precisely
when they all vanish.
\end{proof}

\paragraph{Reproducibility and scope.}
The separate generator \texttt{build\_reciprocal\_reference.py} starts from
the complete shared-reference tables.  It independently substitutes the
rational changes above into all source rows, checks the six resulting
identities exactly, checks that every occurrence of the changed registers
is included, and materializes the full new quadratic systems.  Their
certificates retain \(U_{\mathrm{em}}\) as explicitly unverified.
This one-coordinate improvement applies to the two general interfaces;
no analogous saving for the 56-coordinate ordinary-\(abc\) interface is
asserted here.

\end{document}
